% ============================================================
% Sección 9: Frameworks de agentes (breve)
% ============================================================
\section{Frameworks de agentes}

\begin{frame}{¿Para qué sirve un framework de agentes?}
  \begin{itemize}
    \item Resuelven en común: manejo de historial, parsing de tool calls,
      loops de agente, memoria, orquestación multiagente.
    \item No son estrictamente necesarios (un loop ReAct se puede escribir
      a mano, como veremos en la demo) — pero ahorran trabajo repetitivo
      en proyectos grandes.
    \item Riesgo: añaden \textbf{abstracción y dependencias}; a veces
      conviene empezar simple y adoptar un framework solo cuando se
      necesita.
  \end{itemize}
\end{frame}

\begin{frame}{Panorama de frameworks (2026)}
  \small
  \begin{center}
  \begin{tabular}{@{}p{2.6cm}p{5.6cm}@{}}
    \toprule
    \textbf{Framework} & \textbf{Enfoque principal} \\
    \midrule
    LangChain / LangGraph & Cadenas de LLM + grafos de estados explícitos para agentes complejos \\
    CrewAI & Equipos de agentes con roles definidos (multiagente) \\
    AutoGen / AG2 (Microsoft) & Conversaciones entre múltiples agentes \\
    OpenAI Agents SDK & Agentes ligeros con \emph{handoffs} entre ellos \\
    Claude Agent SDK & El mismo motor de agente detrás de Claude Code, reusable en tus apps \\
    \bottomrule
  \end{tabular}
  \end{center}
\end{frame}

\begin{frame}{¿Framework o loop a mano?}
  \centering
  \resizebox{\textwidth}{!}{%
  \begin{tikzpicture}[node distance=1cm]
    \node[cajaAgente, minimum width=5cm] (simple) {Tarea simple, 1 agente, pocas herramientas};
    \node[cajaHerramienta, below=of simple, minimum width=5cm] (mano) {$\to$ loop ReAct escrito a mano};
    \node[cajaAccento, right=2.5cm of simple, minimum width=5cm] (compleja) {Multiagente, memoria persistente, orquestación};
    \node[cajaLLM, below=of compleja, minimum width=5cm] (fw) {$\to$ framework (LangGraph, CrewAI...)};
    \draw[flecha] (simple) -- (mano);
    \draw[flecha] (compleja) -- (fw);
  \end{tikzpicture}}
  \vfill
  \footnotesize En esta charla construimos el caso simple: entender el loop
  antes de delegarlo a un framework.
\end{frame}

\begin{frame}{Protocolo común: Model Context Protocol (MCP)}
  \begin{itemize}
    \item Estándar abierto (Anthropic, 2024) para conectar un LLM con
      \textbf{herramientas y fuentes de datos} de forma uniforme.
    \item Antes: cada integración era un conector distinto por framework.
    \item Con MCP: un servidor de herramientas puede usarse desde cualquier
      cliente/agente compatible (Claude, editores, otros frameworks).
    \item Ayuda a que el ecosistema de herramientas \textbf{no dependa de
      un solo framework}.
  \end{itemize}
\end{frame}
